\begin{figure*}[t]
\centering
\begin{tikzpicture}[x=1cm,y=1cm,
box/.style={draw=black!45,rounded corners=2pt,align=center,minimum height=1.12cm,font=\small,inner sep=6pt},
flow/.style={-{Latex[length=2mm]},thick,draw=black!65},
sup/.style={flow,dashed,draw=orange!75!black}]
\node[box,fill=blue!7,text width=2.5cm] (build) at (0,0) {Build LiDAR\\endpoints + ray history};
\node[box,fill=blue!7,text width=2.6cm] (shape) at (3.8,0) {Canonical surface\\fixed anchors + 4 children};
\node[box,fill=green!8,text width=2.7cm] (mass) at (7.8,0) {Frozen geometry\\learned F / O / U evidence};
\node[box,fill=green!8,text width=2.4cm] (return) at (11.6,0) {Categorical measure\\interpolated ray return};
\draw[flow] (build)--node[above,font=\scriptsize]{set decoder}(shape);
\draw[flow] (shape)--node[above,font=\scriptsize]{evidence heads}(mass);
\draw[flow] (mass)--node[above,font=\scriptsize]{normalize}(return);
\node[box,fill=orange!8,text width=5.3cm,minimum height=.83cm] (target) at (3.8,-1.72)
{Disjoint target sweeps: set / frame / ray supervision};
\draw[sup] (target.north)--(shape.south);
\draw[sup] (target.north east)--++(0,.25)-|(mass.south);
\node[box,fill=purple!7,text width=3.0cm,minimum height=.83cm] (pose) at (9.1,-1.72)
{Read-only SE(3) pose\\canonical $\leftrightarrow$ world};
\draw[flow] (pose.east)-|(return.south);
\node[font=\scriptsize,align=center,text=black!70] at (4.5,-2.65)
{Physical state $P$ \quad $\xrightarrow{\;\mathrm{stop\mbox{-}gradient}\;}$ \quad visual state $V$ \quad $\xleftarrow{\;\mathrm{image\ loss}\;}$ \quad RGB};
\node[font=\scriptsize,align=center,text=black!70] at (10.8,-2.65)
{Separate static Background};
\end{tikzpicture}
\caption{Overview of Evidential Actor Surfaces. Motion-compensated build observations produce a canonical surface.
Disjoint targets supervise geometry and continuous free/occupied/unknown (F/O/U) evidence in stages.
Frozen evidence weights a categorical ray measure. Rigid pose and appearance have separate state ownership.
Dashed arrows denote supervision.}
\label{fig:eas-overview}
\end{figure*}
